% =====================================================================
\chapter{Identidades trigonométricas}
\label{tema:2}
% =====================================================================

Una identidad trigonométrica es una igualdad entre expresiones trigonométricas que se cumple para todos los valores de la variable en que ambos miembros están definidos. Este tema demuestra las identidades fundamentales a partir de la fórmula del coseno de una diferencia, que a su vez se obtiene de una distancia en la circunferencia unidad, y muestra cómo usarlas para simplificar y verificar expresiones. Se sigue el capítulo~9 de~\cite{osat} (§9.1 a §9.4), equivalente al capítulo~7 de~\cite{osprecalc}; en el \cref{tema:6} se verá una segunda demostración de las fórmulas de la suma, mediante la exponencial compleja, que es la que propone el curso 18.03 del MIT~\cite[§6.4]{mit1803}.

\section{Identidades pitagóricas y verificación de identidades}

\begin{proposicion}[Identidades pitagóricas]\label{prop:pitagoricas}
Para todo $\theta$ en que estén definidas las funciones,
\[
\sin^{2}\theta+\cos^{2}\theta=1,\qquad 1+\tan^{2}\theta=\sec^{2}\theta,\qquad 1+\cot^{2}\theta=\csc^{2}\theta.
\]
\end{proposicion}

\begin{proof}
La primera es el \cref{teo:simetrias}(i). Si $\cos\theta\ne 0$, se divide por $\cos^{2}\theta$ y se obtiene la segunda; si $\sin\theta\ne 0$, se divide por $\sin^{2}\theta$ y se obtiene la tercera.
\end{proof}

Para \emph{verificar} una identidad se transforma uno de sus miembros, mediante identidades ya demostradas y operaciones algebraicas reversibles, hasta obtener el otro. No es válido manipular la igualdad que se quiere probar como si fuera cierta (por ejemplo, elevando ambos miembros al cuadrado), porque así se pueden «demostrar» igualdades falsas~\cite[§9.1]{osat}.

\begin{ejemplo}\label{ej:verificar}
Se verifica $\dfrac{\cos\theta}{1-\sin\theta}=\dfrac{1+\sin\theta}{\cos\theta}$ para $\cos\theta\ne 0$. Esta condición es la que hace que ambos miembros estén definidos: si $\cos\theta\ne0$, entonces $\sin^{2}\theta\ne1$ y, en particular, $1-\sin\theta\ne0$; si $\cos\theta=0$, el miembro derecho no está definido. Multiplicando numerador y denominador del miembro izquierdo por $1+\sin\theta\ne0$:
\[
\frac{\cos\theta\,(1+\sin\theta)}{1-\sin^{2}\theta}=\frac{\cos\theta\,(1+\sin\theta)}{\cos^{2}\theta}=\frac{1+\sin\theta}{\cos\theta}.
\]
\end{ejemplo}

\section{Fórmulas de la suma y de la diferencia}

\begin{teorema}[Coseno de una diferencia]\label{teo:cos-diferencia}
Para todos $\alpha,\beta\in\R$,
\[
\cos(\alpha-\beta)=\cos\alpha\cos\beta+\sin\alpha\sin\beta .
\]
\end{teorema}

\begin{proof}
Se calcula de dos maneras la distancia $d$ entre los puntos $P(\alpha)$ y $P(\beta)$ de la circunferencia unidad. Directamente, usando $\cos^{2}+\sin^{2}=1$,
\[
d^{2}=(\cos\alpha-\cos\beta)^{2}+(\sin\alpha-\sin\beta)^{2}=2-2(\cos\alpha\cos\beta+\sin\alpha\sin\beta).
\]
Por otra parte, el giro de ángulo $-\beta$ alrededor del origen conserva las distancias y las longitudes de arco, y lleva $P(\alpha)$ a $P(\alpha-\beta)$ y $P(\beta)$ a $P(0)=(1,0)$. Por tanto
\[
d^{2}=(\cos(\alpha-\beta)-1)^{2}+\sin^{2}(\alpha-\beta)=2-2\cos(\alpha-\beta).
\]
Igualando las dos expresiones de $d^{2}$ se obtiene el enunciado.
\end{proof}

\begin{corolario}[Fórmulas de la suma y la diferencia]\label{cor:suma}
Para todos $\alpha,\beta\in\R$,
\begin{align*}
\cos(\alpha\pm\beta)&=\cos\alpha\cos\beta\mp\sin\alpha\sin\beta,\\
\sin(\alpha\pm\beta)&=\sin\alpha\cos\beta\pm\cos\alpha\sin\beta,\\
\tan(\alpha\pm\beta)&=\frac{\tan\alpha\pm\tan\beta}{1\mp\tan\alpha\tan\beta}\quad\text{(cuando todas las tangentes están definidas).}
\end{align*}
\end{corolario}

\begin{proof}
Para $\cos(\alpha+\beta)$ se aplica el \cref{teo:cos-diferencia} a $\alpha$ y $-\beta$, con la paridad del \cref{teo:simetrias}(iii). Para el seno, por el \cref{teo:simetrias}(vi),
\[
\sin(\alpha+\beta)=\cos\bigl((\tfrac{\pi}{2}-\alpha)-\beta\bigr)=\cos(\tfrac{\pi}{2}-\alpha)\cos\beta+\sin(\tfrac{\pi}{2}-\alpha)\sin\beta=\sin\alpha\cos\beta+\cos\alpha\sin\beta,
\]
y sustituyendo $\beta$ por $-\beta$ se obtiene $\sin(\alpha-\beta)$. Para la tangente se divide $\sin(\alpha\pm\beta)$ entre $\cos(\alpha\pm\beta)$ y se dividen numerador y denominador por $\cos\alpha\cos\beta\ne0$.
\end{proof}

\begin{ejemplo}\label{ej:cos15}
$\cos 15^{\circ}=\cos(45^{\circ}-30^{\circ})=\frac{\sqrt2}{2}\cdot\frac{\sqrt3}{2}+\frac{\sqrt2}{2}\cdot\frac12=\frac{\sqrt6+\sqrt2}{4}\approx 0{,}9659$.
\end{ejemplo}

\begin{corolario}[Composición de giros]\label{cor:giros}
Sea $R(\alpha)=\begin{pmatrix}\cos\alpha&-\sin\alpha\\ \sin\alpha&\cos\alpha\end{pmatrix}$. Entonces $R(\alpha)R(\beta)=R(\alpha+\beta)$ para todos $\alpha,\beta$.
\end{corolario}

\begin{proof}
La entrada $(1,1)$ del producto es $\cos\alpha\cos\beta-\sin\alpha\sin\beta=\cos(\alpha+\beta)$; la $(2,1)$ es $\sin\alpha\cos\beta+\cos\alpha\sin\beta=\sin(\alpha+\beta)$; la $(1,2)$ es $-\cos\alpha\sin\beta-\sin\alpha\cos\beta=-\sin(\alpha+\beta)$, y la $(2,2)$ es $-\sin\alpha\sin\beta+\cos\alpha\cos\beta=\cos(\alpha+\beta)$.
\end{proof}

La matriz $R(\alpha)$ lleva el vector $(\cos\theta,\sin\theta)$ al vector $(\cos(\theta+\alpha),\sin(\theta+\alpha))$, como se ve aplicando el \cref{cor:suma}: es el giro de ángulo $\alpha$. El \cref{cor:giros} dice que girar $\beta$ y después $\alpha$ equivale a girar $\alpha+\beta$.

\section{Ángulo doble y ángulo mitad}

\begin{proposicion}[Ángulo doble]\label{prop:doble}
Para todo $\alpha$,
\[
\sin 2\alpha=2\sin\alpha\cos\alpha,\qquad \cos 2\alpha=\cos^{2}\alpha-\sin^{2}\alpha=2\cos^{2}\alpha-1=1-2\sin^{2}\alpha,
\]
y $\tan 2\alpha=\dfrac{2\tan\alpha}{1-\tan^{2}\alpha}$ cuando los dos miembros están definidos.
\end{proposicion}

\begin{proof}
Basta tomar $\beta=\alpha$ en el \cref{cor:suma}. Las otras dos formas de $\cos 2\alpha$ se obtienen sustituyendo $\sin^{2}\alpha=1-\cos^{2}\alpha$ o $\cos^{2}\alpha=1-\sin^{2}\alpha$.
\end{proof}

\begin{corolario}[Reducción de potencias y ángulo mitad]\label{cor:mitad}
Para todo $\alpha$,
\[
\cos^{2}\alpha=\frac{1+\cos 2\alpha}{2},\qquad \sin^{2}\alpha=\frac{1-\cos 2\alpha}{2}.
\]
En consecuencia, $\cos\frac{\alpha}{2}=\pm\sqrt{\frac{1+\cos\alpha}{2}}$ y $\sin\frac{\alpha}{2}=\pm\sqrt{\frac{1-\cos\alpha}{2}}$, donde el signo es el de la función en el cuadrante de $\alpha/2$, y $\tan\frac{\alpha}{2}=\frac{1-\cos\alpha}{\sin\alpha}=\frac{\sin\alpha}{1+\cos\alpha}$ cuando los denominadores no se anulan.
\end{corolario}

\begin{proof}
Las dos primeras igualdades se despejan de $\cos2\alpha=2\cos^{2}\alpha-1=1-2\sin^{2}\alpha$. Sustituyendo $\alpha$ por $\alpha/2$ y tomando raíces se obtienen las fórmulas del ángulo mitad; la raíz cuadrada solo determina el valor absoluto, y el signo depende del cuadrante. Para la tangente, $\frac{1-\cos\alpha}{\sin\alpha}=\frac{2\sin^{2}(\alpha/2)}{2\sin(\alpha/2)\cos(\alpha/2)}=\tan\frac{\alpha}{2}$; la otra forma se obtiene igual con $1+\cos\alpha=2\cos^{2}(\alpha/2)$.
\end{proof}

\begin{ejemplo}\label{ej:mitad}
$\sin\frac{\pi}{12}=+\sqrt{\frac{1-\cos(\pi/6)}{2}}=\sqrt{\frac{2-\sqrt3}{4}}=\frac{\sqrt{2-\sqrt3}}{2}\approx 0{,}2588$, con signo positivo porque $\pi/12$ está en el primer cuadrante.
\end{ejemplo}

\section{Fórmulas de producto a suma y de suma a producto}

\begin{proposicion}[Producto a suma]\label{prop:producto-suma}
Para todos $\alpha,\beta$,
\begin{align*}
\cos\alpha\cos\beta&=\tfrac12\bigl[\cos(\alpha-\beta)+\cos(\alpha+\beta)\bigr],\\
\sin\alpha\sin\beta&=\tfrac12\bigl[\cos(\alpha-\beta)-\cos(\alpha+\beta)\bigr],\\
\sin\alpha\cos\beta&=\tfrac12\bigl[\sin(\alpha+\beta)+\sin(\alpha-\beta)\bigr].
\end{align*}
\end{proposicion}

\begin{proof}
Se suman o restan las fórmulas del \cref{cor:suma}: $\cos(\alpha-\beta)+\cos(\alpha+\beta)=2\cos\alpha\cos\beta$; $\cos(\alpha-\beta)-\cos(\alpha+\beta)=2\sin\alpha\sin\beta$; $\sin(\alpha+\beta)+\sin(\alpha-\beta)=2\sin\alpha\cos\beta$.
\end{proof}

\begin{corolario}[Suma a producto]\label{cor:suma-producto}
Para todos $x,y$,
\begin{align*}
\sin x+\sin y&=2\sin\tfrac{x+y}{2}\cos\tfrac{x-y}{2}, & \sin x-\sin y&=2\cos\tfrac{x+y}{2}\sin\tfrac{x-y}{2},\\
\cos x+\cos y&=2\cos\tfrac{x+y}{2}\cos\tfrac{x-y}{2}, & \cos x-\cos y&=-2\sin\tfrac{x+y}{2}\sin\tfrac{x-y}{2}.
\end{align*}
\end{corolario}

\begin{proof}
En la \cref{prop:producto-suma} se toma $\alpha=\frac{x+y}{2}$ y $\beta=\frac{x-y}{2}$, de modo que $\alpha+\beta=x$ y $\alpha-\beta=y$. Por ejemplo, la tercera línea de la proposición da $2\sin\frac{x+y}{2}\cos\frac{x-y}{2}=\sin x+\sin y$. Las demás son análogas; para la segunda se aplica la tercera línea con $\alpha=\frac{x-y}{2}$ y $\beta=\frac{x+y}{2}$, y se usa que el seno es impar.
\end{proof}

\begin{observacion}[Errores frecuentes]
(a) Las funciones trigonométricas no son lineales: $\sin(\alpha+\beta)\ne\sin\alpha+\sin\beta$ y $\sin 2\alpha\ne 2\sin\alpha$ en general (para $\alpha=\beta=\pi/4$ el primer miembro vale $1$ y el segundo $\sqrt2$). (b) En las fórmulas del ángulo mitad, el signo de la raíz no es opcional: lo fija el cuadrante de $\alpha/2$, no el de $\alpha$. (c) Al verificar una identidad hay que indicar dónde están definidos ambos miembros.
\end{observacion}

\begin{observacion}[Conexión con el curso]
El \cref{cor:giros} es exactamente el lema de rotaciones con el que se demuestra que RoPE solo depende de la posición relativa (módulo~14, tema~4). Las fórmulas de la suma explican también la propiedad que Vaswani \emph{et al.}~\cite[§3.5]{vaswani} señalan para la codificación sinusoidal: para cada frecuencia $\omega$, el par $(\sin\omega(p+k),\cos\omega(p+k))$ se obtiene del par $(\sin\omega p,\cos\omega p)$ multiplicando por una matriz que depende de $k$ pero no de $p$ (\cref{eje:4-pe-lineal}).
\end{observacion}

\begin{lecturas}
\item OpenStax, \emph{Algebra and Trigonometry}~\cite{osat}, §9.1 a §9.4 (identidades, suma y diferencia, ángulo doble y mitad, producto y suma). La numeración equivale a §7.1 a §7.4 de \emph{Precalculus}~\cite{osprecalc}.
\item MIT 18.03, \emph{Supplementary Notes}~\cite{mit1803}, §6.4 «Multiplication», donde las fórmulas de la suma se deducen de la exponencial compleja.
\end{lecturas}

\section*{Ejercicios del tema 2}
\addcontentsline{toc}{section}{Ejercicios del tema 2}

\begin{ejercicio}{1}{eje:2-vf}
Justifique si son verdaderas o falsas para todo $\alpha,\beta$: (a) $\cos(\alpha+\beta)=\cos\alpha+\cos\beta$; (b) $\sin2\alpha=2\sin\alpha\cos\alpha$; (c) $\cos2\alpha=1-2\cos^{2}\alpha$; (d) $\tan(\alpha+\pi)=\tan\alpha$ donde esté definida.
\end{ejercicio}

\begin{ejercicio}{1}{eje:2-signo-mitad}
Sea $\pi<\alpha<2\pi$. ¿Qué signo hay que tomar en $\cos\frac{\alpha}{2}=\pm\sqrt{(1+\cos\alpha)/2}$ y en $\sin\frac{\alpha}{2}=\pm\sqrt{(1-\cos\alpha)/2}$? Justifique la respuesta.
\end{ejercicio}

\begin{ejercicio}{2}{eje:2-exactos}
Calcule el valor exacto de $\sin 75^{\circ}$, $\cos 105^{\circ}$ y $\tan 105^{\circ}$ con las fórmulas de la suma.
\fuente{adaptado de OpenStax, \emph{Algebra and Trigonometry}~\cite{osat}, §9.2, ejercicios de la sección sobre valores exactos.}
\end{ejercicio}

\begin{ejercicio}{2}{eje:2-cuadrantes}
Se sabe que $\sin\alpha=3/5$ con $\alpha$ en el primer cuadrante y $\cos\beta=-5/13$ con $\beta$ en el segundo. Calcule $\sin(\alpha+\beta)$, $\cos(\alpha+\beta)$ y $\tan(\alpha+\beta)$, y determine el cuadrante de $\alpha+\beta$.
\fuente{adaptado de OpenStax, \emph{Algebra and Trigonometry}~\cite{osat}, §9.2, ejercicios de la sección con datos en distintos cuadrantes.}
\end{ejercicio}

\begin{ejercicio}{2}{eje:2-doble}
Se sabe que $\cos\theta=-3/5$ y que $\theta$ está en el tercer cuadrante. Calcule $\sin2\theta$, $\cos2\theta$ y $\tan2\theta$.
\fuente{adaptado de OpenStax, \emph{Algebra and Trigonometry}~\cite{osat}, §9.3, ejercicios de la sección sobre el ángulo doble.}
\end{ejercicio}

\begin{ejercicio}{2}{eje:2-pi8}
Calcule el valor exacto de $\sin\frac{\pi}{8}$, $\cos\frac{\pi}{8}$ y $\tan\frac{\pi}{8}$, simplificando este último hasta la forma $a+b\sqrt2$.
\end{ejercicio}

\begin{ejercicio}{2}{eje:2-producto}
Escriba $\cos3x\cos x$ y $\sin4x\cos2x$ como sumas, y $\sin5x-\sin3x$ y $\cos x+\cos3x$ como productos.
\fuente{adaptado de OpenStax, \emph{Algebra and Trigonometry}~\cite{osat}, §9.4, ejercicios de la sección sobre producto a suma y suma a producto.}
\end{ejercicio}

\begin{ejercicio}{2}{eje:2-potencias}
Exprese $\sin^{4}x$ como combinación lineal de $1$, $\cos2x$ y $\cos4x$.
\end{ejercicio}

\begin{ejercicio}{2}{eje:2-verificar}
Verifique las identidades (a) $\dfrac{1-\cos2x}{\sin2x}=\tan x$ y (b) $\dfrac{\sin x}{1+\cos x}+\dfrac{1+\cos x}{\sin x}=2\csc x$, e indique en qué valores de $x$ son válidas.
\end{ejercicio}

\begin{ejercicio}{2}{eje:2-triple}
Demuestre las fórmulas del ángulo triple $\cos3x=4\cos^{3}x-3\cos x$ y $\sin3x=3\sin x-4\sin^{3}x$.
\end{ejercicio}

\begin{ejercicio}{3}{eje:2-cos36}
Sea $\theta=\pi/5$. Usando que $3\theta=\pi-2\theta$, demuestre que $\cos\theta$ es raíz de $4c^{2}-2c-1=0$ y deduzca que $\cos\frac{\pi}{5}=\frac{1+\sqrt5}{4}$.
\end{ejercicio}

\begin{ejercicio}{3}{eje:2-suma-senos}
Demuestre que, si $\sin(x/2)\ne 0$, para todo $n\ge1$
\[
\sum_{k=1}^{n}\sin kx=\frac{\sin\frac{nx}{2}\,\sin\frac{(n+1)x}{2}}{\sin\frac{x}{2}}.
\]
Indicación: multiplique la suma por $2\sin(x/2)$ y use la \cref{prop:producto-suma} para obtener una suma telescópica.
\end{ejercicio}

\begin{ejercicio}{3}{eje:2-tan-triangulo}
Sean $\alpha,\beta,\gamma$ los ángulos de un triángulo no rectángulo. Demuestre que $\tan\alpha+\tan\beta+\tan\gamma=\tan\alpha\tan\beta\tan\gamma$.
\end{ejercicio}

\begin{ejercicio}{3}{eje:2-chebyshev}
Demuestre por inducción que, para cada $n\ge1$, existe un polinomio $T_n$ de grado $n$ y coeficiente principal $2^{n-1}$ tal que $\cos nx=T_n(\cos x)$ para todo $x$. Indicación: pruebe primero que $\cos(n+1)x+\cos(n-1)x=2\cos x\cos nx$.
\end{ejercicio}

\begin{ejercicio}{3}{eje:2-potencia-giro}
Con la notación del \cref{cor:giros}, demuestre por inducción que $R(\alpha)^{n}=R(n\alpha)$ para todo entero $n\ge0$, que $R(\alpha)$ es invertible con $R(\alpha)^{-1}=R(-\alpha)=R(\alpha)\T$, y deduzca que la igualdad $R(\alpha)^{n}=R(n\alpha)$ vale para todo $n\in\Z$.
\end{ejercicio}
